\section{\label{I_3_B}Segmentation} 

Dans le cas des archives audiovisuelles, une indexation fine ne peut s’effectuer à l’échelle du fichier brut dans sa totalité, en particulier lorsqu’il s’agit de rushes contenant une succession importante de séquences non montées. Afin de rendre le contenu de telles archives exploitable au sein d’un modèle de données cohérent, il est indispensable d’analyser finement le flux continu de l’image via une tâche de segmentation. Cette opération est à l’origine de questionnements autour de l’intelligibilité du document audiovisuel :  

\begin{quote} 
	
	"Quels sont les mécanismes qui permettent d’in-terpréter cette longue séquence d’images [...] en une séquence d’actions dramatiques, reliées entre elles par des liens de causalité, et qui nous permettent d’en comprendre le sens (l’histoire) ? "\footcite{zotero-item-644} 
	
\end{quote} 

L’intelligibilité du document audiovisuel passerait ainsi par une interprétation, via une segmentation de l’enchaînement des images afin d’en extraire le sens. La segmentation revêt, dans le cadre de l’indexation d’archives audiovisuelles, une dimension essentiellement temporelle, dans l'optique de découper un flux vidéo selon certains critères. Il s'agit d'un découpage conceptuel et non d'un montage technique qui modifie le sens de l’image animée. La segmentation s'imposerait comme la condition d'une indexation fine puisqu'en apposant à des séquences des balises temporelles, elle permettrait de rendre possible "la production de descriptions structurées". \footcite{stockinger2011} Ces balises temporelles délimitent des ensembles homogènes au sein d'une archive audiovisuelle, à la manière de chapitres dans un ouvrage textuel \footcite{stockinger2011}. Si la structuration de certains types de contenus audiovisuels tels que les journaux télévisés, le sport, la publicité, sont "relativement stable et explicite de leur structure", et peuvent alors facilement être segmentés, les archives audiovisuelles regroupent des corpus divers et variés à la structure aléatoire qui rendent complexe cette tâche \footcite{carrive2014} 

La segmentation s’applique particulièrement bien au corpus d’archives audiovisuelles composé de rushes documentaires, que nous avons décrit. En effet, au sein du logiciel prévu par Skinsoft, la segmentation est pratiquée par l’institution patrimoniale préservant ces archives, qui associe des repères temporels à des extraits numérisés et issus de rushes filmés plus larges. Cette pratique permet d’extraire des rushes certaines séquences jugées importantes à indexer de manière précise. Ces séquences sont ensuite rendues intelligibles par l’indexation de leur contenu précis et deviennent requêtables de manière autonome au sein du modèle de données.  

La segmentation répond donc à un objectif fixé par l’indexation, celui d’agir comme “un filtre de pertinence" sur les éléments à indexer \footcite{stockinger2011}. Sa mise en place suppose d’en définir la granularité (analyse au niveau du plan, du changement de plan, de scènes, de séquences, etc), afin d’isoler les éléments d’analyse \footcite{stockinger2011}. Ce travail d’analyse s’articule autour de trois types de segmentation, conditionnée par la nature du médium cinématographique. En effet, la segmentation temporelle est non seulement conditionnée par l’image mais également par le son, composant essentiel des documents audiovisuels \footcite{carrive2014}. La segmentation temporelle suppose également une segmentation spatiale, induisant l'analyse de la situation des éléments dans l'image \footcite{zhou2022}. Enfin, une segmentation sémantique intervient, apposant une description aux éléments extraits, permettant de produire un découpage logique \footcite{zhou2022}. Les catégories sémantiques sont prédéfinies par l'utilisateur et peuvent être de plusieurs types conformément à la nature de l’image animée. Ces catégories peuvent être d’ordre technique, relationnelles, de contenu, de mouvement, etc.  

Parmi ces trois types de segmentation, seule la segmentation temporelle est spécifique à l'image animée, puisqu'elle analyse la différence temporelle et de mouvement entre les plans ainsi que les relations qu’ils entretiennent, afin de pouvoir en extraire des séquences distinctes. Les segmentation spatiales et sémantiques sont, en revanche, également applicables à l'image fixe. 

Ce processus de segmentation s’inscrit directement dans l’architecture du modèle FRBR. Les segments temporels créés sont décrits et indexés comme des nouvelles entités et rattachés aux entités pertinentes du modèle. C’est dans le contexte de l’indexation de grands volumes d’archives et spécifiquement de corpus complexes que peut être envisagée la mise en place d’une automatisation de la segmentation, puis de l‘indexation. Selon Attard, la mise en place d’une segmentation automatisée permettrait “d’accélérer le traitement de grands volumes vidéo tout en garantissant la cohérence du processus” (traduction : automated segmentation accelerates the processing of large volumes of video content and ensures consistency and accuracy in the segmentation process." ) \footcite{attard2025}. 

Toutefois, l’efficacité promise par l’automatisation de la segmentation et de l’indexation ne pourra être mise en place qu’à la condition selon laquelle les algorithmes d’apprentissage automatique puissent produire ces tâches selon des règles métiers précises et un modèle de données explicite. En tant qu’éditeur de logiciel de gestion des collections, Skinsoft ambitionne de résoudre ces complexités afin d’intégrer ces nouvelles fonctionnalités au sein de l’espace de travail S-Movies. Dès lors, il convient d’examiner les directions adoptées par l’éditeur afin de proposer une automatisation innovante conditionnée par une structure de données complexe.  